\section{Introduction}

Next Best View computation~(NBV) is a long-standing problem in robotics~\cite{connolly-father-paper, yamauchi-97}, which consists in identifying the next most informative sensor position(s) for reconstructing a 3D object or scene efficiently and accurately. Typically, a position is evaluated on how much it can increase the total coverage of the scene surface. 
Few methods have relied on Deep Learning (DL) for the NBV problem, even though DL can provide useful geometric prior to obtain a better prediction of the surface coverage~\cite{zeng-icirs20-pcnbv, Mendoza-2019, vasquez-21-3DCNN}. 
Like most current methods, we consider NBV prediction from a depth sensor. Existing methods based on a depth sensor rely either on a volumetric or on a surface-based representation of the scene geometry. Volumetric mapping-based methods can compute collision efficiently, which is practical for path planning in real case scenarios~\cite{potthast-14-probabilistic-framework, vasquez-14-volumetric-nbv, vasquez-17-view-state-planning, kriegel-12-next-best-scan, bissmarck-15-efficient-nbv, daudelin-17-adaptable-probabilistic}. However, they typically rely on voxels or a global embedding~\cite{hepp-corr18-learntoscore, charrow-15-information, song-17-online-inspection, song-20-online-coverage, wang-20-efficient-autonomous, cieslewski-17-ICIRS} for the scene, which results in poor accuracy in reconstruction and poor performance in NBV selection for complex 3D objects. On the contrary, surface mapping-based methods that process directly a dense point cloud of the surface as gathered by the depth sensor are efficient for NBV prediction with high-detailed geometry. They are however limited to very specific cases, generally a single, small-scale, isolated object with the camera constrained to stay on a sphere centered on the object~\cite{lee-tra20-mechanical-parts,vasquez-17-view-state-planning, delmerico-18-acomparison, chen-smc05-visionsensorplanning,kriegel-11-surface-based, kriegel-12-next-best-scan,kriegel-12-next-best-scan,zeng-icirs20-pcnbv, Mendoza-2019, vasquez-21-3DCNN}. Thus, they cannot be applied to the exploration of 3D scenes. 


% They are however limited to very specific cases~\cite{lee-tra20-mechanical-parts}, generally a single, small-scale~\cite{vasquez-17-view-state-planning, delmerico-18-acomparison, chen-05-visionsensorplanning, kriegel-11-surface-based, kriegel-12-next-best-scan}, isolated object~\cite{kriegel-12-next-best-scan} floating in the air with the camera constrained to stay on a sphere centered on the object~\cite{zeng-icirs20-pcnbv, Mendoza-2019, vasquez-21-3DCNN}. Thus, they cannot be applied to the exploration of 3D scenes. 


\begin{figure}%
    \centering
    \begin{tabular}{cccc}
      \includegraphics[width=3.15cm, trim={0 0 3cm 0}, clip]{images/colosseum1.png} & % width=0.45\linewidth
      \includegraphics[width=3.15cm, trim={3cm 0 3cm 0}, clip]{images/pc_compressed/colosseum_recovered_2.png} &
      \includegraphics[width=3.15cm, trim={0 0 0 0}, clip]{images/compressed/neuschwanstein.png} & % width=0.45\linewidth
      \includegraphics[width=2.8cm, trim={6cm 0 3cm 3cm}, clip]{images/pc_compressed/bavaria.png}\\
    \end{tabular}
    \caption{{\bf Our Next Best View~(NBV) method \method  can handle unknown large-scale 3D scenes} to produce accurate 3D reconstruction inside a given 3D bounding box. Here, we call \method iteratively within a naive path planning algorithm to compute a complete camera trajectory avoiding collisions and obtain a complete 3D model. Despite being trained only on ShapeNet 3D models, it generalizes to complex scenes as shown above. (input 3D model courtesy of Brian Trepanier, under CC License. Downloaded from Sketchfab.)}%
    \label{fig:teaser}%
\end{figure}


As shown in Figure~\ref{fig:teaser}, we introduce a volumetric DL method that efficiently identifies NBVs for unknown large-scale 3D scenes in which the camera can move freely.  Instead of representing the scene with a single global embedding, we choose to use a grid of local point clouds, which scales much better to large and complex 3D scenes. We show how to learn to predict the visibility of unseen 3D points in all directions given an estimate of the 3D scene geometry. We can then integrate these visibilities in the direction of any camera by using a Monte Carlo integration approach, which allows us to optimize the camera pose to find the next most informative views.  We call our method \method, for Surface Coverage Optimization in uNknown Environments.

In this respect, we introduce a theoretical framework to translate the optimization of surface coverage gain, a surface metric on manifolds that represents the ability of a camera to increase the visible area of the surface, into an optimization problem on volumetric integrals. Such a formalism allows us to use a volumetric mapping of geometry, which is convenient not only to scale the model to exploration tasks and scene reconstruction, but also to make probabilistic predictions about geometry.

In particular, given a partial point cloud gathered by a depth sensor, our model learns to make predictions with virtually infinite resolution about the occupancy probability in the scene volume by learning a deep implicit function~\cite{lars-18-occupancy_network, xu-19-disn, mildenhall-20-nerf, yariv-2020-idr, yariv-2021-volsdf, oechsle-21-unisurf}. Such predictions scale to very large point clouds since they depend only on neighborhood geometry. Then, our model leverages a self-attention mechanism~\cite{ashish-2017-attention} to predict occlusions and compute informative functions mapped on a continuous sphere that represent visibility likelihood of points in all directions. The occupancy probability field is finally used as a basis to sample points and compute Monte Carlo integrals of visibility scores.

Since NBV learning-based methods are mostly limited to single, small-scale, centered object reconstruction in literature, we first compare the performance of our model to the state of the art on the ShapeNet dataset~\cite{shapenet2015}, following the protocol introduced in \cite{zeng-icirs20-pcnbv}. While our method was designed to handle more general frameworks such as 3D scene reconstruction and continuous cameras poses in the scene, it outperforms the state of the art for dense reconstruction of objects when the camera is constrained to stay on a sphere centered on the object. We then conduct experiments in large 3D environments using a simple planning algorithm that builds a camera trajectory online by iteratively selecting NBVs with \method.  Since, to the best of our knowledge, we propose the first supervised Deep Learning method for such free 6D motion of the camera, we created a dataset made of several large-scale scenes under the CC License for quantitative evaluation. We will make our code and this dataset available for allowing comparison of future methods with \method on our project webpage: \url{https://github.com/Anttwo/SCONE}.


\vincentrmk{What about the code?}


\antoinermk{Indeed we are the first, for NBV selection of structures with detailed geometry (with free 6D motion, it's true). Maybe we should insist on this point; there are some works on path planning with coarse geometry, but their focus is not on high quality of reconstruction. There are also some RL approaches too but it's really different and in this particular field they all look quite the same in my opinion ahah. I should check once again the SOTA to be sure}
